\documentclass[a4paper]{article}


\usepackage{graphicx}
\graphicspath{}
\DeclareGraphicsExtensions{.pdf,.png,.jpg}
\usepackage{caption}

\usepackage{booktabs}
\usepackage{array}
\usepackage{multirow}
\usepackage{longtable}
\usepackage{xcolor}
\usepackage{tcolorbox}
\usepackage{titlesec}

\usepackage{tikz}  
\usetikzlibrary{graphs}
\usetikzlibrary{positioning,arrows}
\usetikzlibrary{trees}

\usepackage{fancyhdr}

\usepackage[unicode=true, colorlinks=true, linkcolor=blue, urlcolor=blue]{hyperref}

\usepackage{header}
\usepackage{enumitem}
\setlist{nolistsep, itemsep=0.1cm,parsep=0pt}

%Russian-specific packages
%--------------------------------------
{\usepackage[T2A]{fontenc}
	\usepackage[utf8]{inputenc}
	\usepackage[russian]{babel}
	\usepackage{datetime}}


%--------------------------------------

%Hyphenation rules
%--------------------------------------
\usepackage{hyphenat}
\hyphenation{ма-те-ма-ти-ка вос-ста-нав-ли-вать}
\renewcommand{\dateseparator}{.}
%--------------------------------------






\pagestyle{fancy}
\fancyhf{}
\fancyhead[L]{\small Экономическая статистика}
\fancyhead[R]{\small Справочные материалы}
\fancyfoot[C]{\thepage}


\begin{document}
	
	
\tableofcontents

\section{Описательная статистика}
% =============================================================

Описательная статистика позволяет резюмировать набор данных с помощью числовых мер.
Ключевые категории:

\begin{itemize}[noitemsep]
	\item \textbf{Характеристики центра распределения} - описывают «типичное» значение
	(среднее, медиана, мода).
	\item \textbf{Порядковые статистики} - показывают место наблюдения в распределении
	(квантили, перцентили).
	\item \textbf{Меры формы} - описывают форму распределения
	(асимметрия, эксцесс).
\end{itemize}

\noindent\textbf{Обозначения:} 

\textbf{$n$ }- объём выборки; 

\textbf{$N$} - объём генеральной
совокупности; 

\textbf{$x_i$} - $i$-е наблюдение; 

\textbf{$f_i$} - частота $i$-го значения; нижний индекс~$0$ - базисный период, $t$ - отчётный период.


% =============================================================
\section{Средние величины, меры центра распределения}
% =============================================================

% -----------------------------------------------------------
\subsection{Среднее арифметическое}
% -----------------------------------------------------------


	\textbf{Определение.} Среднее арифметическое - сумма всех значений, делённая
	на их количество.


\subsubsection{Простое среднее арифметическое}

\noindent\textbf{Для выборки:}
\begin{equation}
	\bar{x} = \frac{\displaystyle\sum_{i=1}^{n} x_i}{n}
\end{equation}

\noindent\textbf{Для генеральной совокупности:}
\begin{equation}
	\mu = \frac{\displaystyle\sum_{i=1}^{N} x_i}{N}
\end{equation}

\subsubsection{Взвешенное среднее арифметическое}

Используется, когда отдельные значения встречаются с различной частотой $f_i$:

\begin{equation}
	\bar{x}_w = \frac{\displaystyle\sum_{i=1}^{k} x_i f_i}{\displaystyle\sum_{i=1}^{k} f_i}
	= \frac{\displaystyle\sum_{i=1}^{k} x_i f_i}{n}
\end{equation}


	\textbf{Пример.} Цены на товар в трёх регионах: 100, 120, 130~р.
	Объёмы продаж: 200, 500, 300~единиц.
	\[
	\bar{x}_w = \frac{100\cdot200 + 120\cdot500 + 130\cdot300}{200+500+300}
	= \frac{20000+60000+39000}{1000} = 119 \text{ р.}
	\]
	Простое среднее: $(100+120+130)/3 \approx 116{,}7$~р. - менее точный результат,
	так как не учитывает объёмы продаж.


\subsubsection{Свойства среднего арифметического}

\begin{enumerate}[noitemsep]
	\item $\sum (x_i - \bar{x}) = 0$ - сумма отклонений от среднего равна нулю.
	\item $\sum (x_i - \bar{x})^2 \to \min$ - сумма квадратов отклонений
	минимальна относительно любой другой константы.
	\item Чувствительно к выбросам (аномальным наблюдениям).
	\item Применяется только к количественным данным.
\end{enumerate}

% -----------------------------------------------------------
\subsection{Медиана}
% -----------------------------------------------------------


	\textbf{Определение.} Медиана $(\bar{x}$ или $\text{Me})$ - значение,
	делящее упорядоченный ряд данных на две равные части: 50\% наблюдений ниже
	медианы и 50\% выше.


\noindent\textbf{Несгруппированные данные.}
Данные упорядочиваются по возрастанию, тогда:
\[
\bar{x} =
\begin{cases}
	x_{(n+1)/2}, & n \text{ - нечётное,}\\[4pt]
	\dfrac{x_{n/2} + x_{n/2+1}}{2}, & n \text{ - чётное.}
\end{cases}
\]

\noindent\textbf{Сгруппированные данные.}
Находится медианный класс - первый класс с накопленной частотой $\geq n/2$:

\begin{equation}
	\bar{x} = x_0 + h \cdot \frac{\dfrac{n}{2} - S_{Me-1}}{f_{Me}} 
\end{equation}

где $x_0$ - нижняя граница медианного класса;
$ S_{ПС-1}$ - накопленная частота до медианного класса;
${f_{ПС}}$ - частота медианного класса; $h$ - ширина интервала.


	\textbf{Когда использовать медиану вместо среднего?}
	Медиана предпочтительна при наличии выбросов или при скошенных распределениях: доходы населения, цены на жильё, время ожидания.


% -----------------------------------------------------------
\subsection{Мода}
% -----------------------------------------------------------


	\textbf{Определение.} Мода $(\text{Mo})$ - значение, которое
	встречается наиболее часто.


\noindent\textbf{Для сгруппированных данных и равными интервалами:}
\begin{equation}
	\text{Mo} = x_0 + h_{Mo} \cdot \frac{f_m - f_{m-1}}{(f_m - f_{m-1}) + (f_m - f_{m+1})} 
\end{equation}

где $f_m$ - частота модального класса;
$f_{m-1}$, $f_{m+1}$ - частоты соседних классов.

\noindent\textbf{Для неравных интервалов:}

При работе с неравными интервалами мы сначала расчитываем относительную частоту для каждого интервала $ \rho_i = \frac{f_i}{h_i} $. И определяется модальный интервал по столбцу с относительными частотами $ \max_{\rho}(\text{Mo}_i) $

\begin{equation}
	\text{Mo} = x_0 + h_{Mo} \cdot \frac{\rho_m - \rho_{m-1}}{(\rho_m - \rho_{m-1}) + (\rho_m - \rho_{m+1})} 
\end{equation}

\begin{itemize}[noitemsep]
	\item \textbf{Унимодальное} распределение: одна мода.
	\item \textbf{Бимодальное}: две равнозначные моды.
	\item \textbf{Мультимодальное}: несколько мод.
	\item Если все значения различны - мода отсутствует.
\end{itemize}



% -----------------------------------------------------------
\subsection{Среднее геометрическое}
% -----------------------------------------------------------


	\textbf{Определение.} Среднее геометрическое - корень степени $n$ из произведения
	$n$ положительных значений.


\begin{equation}
	\bar{x}_{геом} = \sqrt[n]{\prod_{i=1}^{n} x_i} = \left(\prod_{i=1}^{n} x_i\right)^{1/n}
\end{equation}

\noindent\textbf{Логарифмический эквивалент:}
\[
\ln \bar{x}_{геом} = \frac{1}{n}\sum_{i=1}^{n} \ln x_i
\]


	\textbf{Пример: среднегодовой темп роста.}
	Выручка компании росла: $+8\%$, $+12\%$, $+3\%$, $+15\%$ за 4~года.
	Индексы роста: $1{,}08;\; 1{,}12;\; 1{,}03;\; 1{,}15$.
	\[
	\bar{x}_{геом} = (1{,}08 \times 1{,}12 \times 1{,}03 \times 1{,}15)^{0{,}25}
	= (1{,}4321)^{0{,}25} \approx 1{,}0941
	\]
	Среднегодовой темп роста $\approx 9{,}41\%$.



	\textbf{Область применения:} темпы роста, доходность инвестиций, индексы изменения
	цен за несколько периодов. Все значения $x_i$ должны быть строго положительными.


% -----------------------------------------------------------
\subsection{Среднее гармоническое}
% -----------------------------------------------------------


	\textbf{Определение.} Среднее гармоническое - величина, обратная среднему
	арифметическому обратных значений.


\begin{equation}
	\bar{x}_{гарм} = \frac{n}{\displaystyle\sum_{i=1}^{n} \dfrac{1}{x_i}}
\end{equation}

\noindent\textbf{Взвешенное среднее гармоническое:}
\begin{equation}
	\bar{x}_{геом}^w = \frac{\displaystyle\sum_{i=1}^{k} w_i}{\displaystyle\sum_{i=1}^{k} \dfrac{w_i}{x_i}}
\end{equation}


	\textbf{Пример: средняя скорость.}
	Автомобиль проехал три равных участка со скоростями 60, 80 и 100~км/ч.
	\[
	\bar{x}_{геом} = \frac{3}{\dfrac{1}{60}+\dfrac{1}{80}+\dfrac{1}{100}}
	= \frac{3}{0{,}01667+0{,}01250+0{,}01000}
	= \frac{3}{0{,}03917} \approx 76{,}6 \text{ км/ч}
	\]



	\textbf{Область применения:} средние цены при фиксированных суммах затрат,
	производительность труда, скорость при одинаковых расстояниях.


% -----------------------------------------------------------
\subsection{Среднее квадратическое}
% -----------------------------------------------------------

\begin{equation}
	\bar{x}_{\text{кв}} = \sqrt{\frac{\displaystyle\sum_{i=1}^{n} x_i^2}{n}}
\end{equation}

Применяется в физике, электротехнике и при вычислении стандартного отклонения, когда исходные значения положительные и отрицательные .

% -----------------------------------------------------------
\subsection{Средняя хронологическая}
% -----------------------------------------------------------


	\textbf{Определение.} Средняя хронологическая - средняя из значений показателя,
	зафиксированных в последовательные моменты времени. Применяется для нахождения
	среднего уровня \emph{моментного} временного ряда (например, остатки товаров
	на складе, численность персонала, объём вкладов на конкретные даты).


\subsubsection{Равноотстоящие моменты времени}

Если интервалы между датами одинаковы, то:

\begin{equation}
	\bar{x}_{\text{хр}} =
	\frac{\dfrac{x_1}{2} + x_2 + x_3 + \cdots + x_{n-1} + \dfrac{x_n}{2}}{n - 1}
\end{equation}

Крайние члены ряда берутся с коэффициентом $\tfrac{1}{2}$, все остальные - с
коэффициентом $1$.


	\textbf{Пример.} Остатки товаров на складе (тыс. ед.) на начало каждого месяца:
	
	\begin{center}
		\begin{tabular}{ccccccc}
			\toprule
			Дата & 1 янв & 1 фев & 1 мар & 1 апр & 1 май & 1 июн \\
			\midrule
			$x_i$ & 120 & 140 & 100 & 160 & 130 & 110 \\
			\bottomrule
		\end{tabular}
	\end{center}
	
	$n = 6$ моментов, интервалы равны (1 месяц). Применяем формулу:
	
	\[
	\bar{x}_{\text{хр}}
	= \frac{\dfrac{120}{2} + 140 + 100 + 160 + 130 + \dfrac{110}{2}}{6 - 1}
	= \frac{60 + 140 + 100 + 160 + 130 + 55}{5}
	= \frac{645}{5} = 129 \text{ тыс. ед.}
	\]
	
	\textbf{Интерпретация:} средний остаток товаров на складе за данный период
	составил 129~тыс. единиц.


\subsubsection{Неравноотстоящие моменты времени}

Если интервалы между датами различны (длины $t_1, t_2, \ldots, t_{n-1}$),
применяется взвешенная формула - каждый средний уровень на отрезке
взвешивается на длину соответствующего интервала:

\begin{equation}
	\bar{x}_{\text{хр}} =
	\frac{\displaystyle\sum_{j=1}^{n-1} \frac{x_j + x_{j+1}}{2} \cdot t_j}
	{\displaystyle\sum_{j=1}^{n-1} t_j}
\end{equation}

где $t_j$ - длина $j$-го интервала (в днях, месяцах или любых других единицах).


	\textbf{Пример.} Численность сотрудников компании на ключевые даты:
	
	\begin{center}
		\renewcommand{\arraystretch}{1.3}
		\begin{tabular}{cccc}
			\toprule
			Дата & Числ., чел. & Следующая дата & Интервал $t_j$, дн. \\
			\midrule
			1 янв  & 200 & 1 мар  & 59 \\
			1 мар  & 220 & 1 июл  & 122 \\
			1 июл  & 180 & 1 окт  & 92 \\
			1 окт  & 210 & 1 янв  & 92 \\
			\bottomrule
		\end{tabular}
	\end{center}
	
	\begin{align*}
		\bar{x}_{\text{хр}} &=
		\frac{\dfrac{200+220}{2}\cdot59
			+ \dfrac{220+180}{2}\cdot122
			+ \dfrac{180+210}{2}\cdot92
			+ \dfrac{210+200}{2}\cdot92}
		{59 + 122 + 92 + 92} \\[8pt]
		&= \frac{210\cdot59 + 200\cdot122 + 195\cdot92 + 205\cdot92}{365} \\[4pt]
		&= \frac{12390 + 24400 + 17940 + 18860}{365}
		= \frac{73590}{365} \approx 201{,}6 \text{ чел.}
	\end{align*}
	
	\textbf{Интерпретация:} средняя за год численность персонала составила около
	202~человек.


\subsubsection{Отличие от интервального ряда}

\begin{center}
	\renewcommand{\arraystretch}{1.4}
	\begin{tabular}{p{3.8cm} p{4.5cm} p{4.5cm}}
		\toprule
		& \textbf{Моментный ряд} & \textbf{Интервальный ряд} \\
		\midrule
		Что фиксируется & Уровень на дату & Итог за период \\
		Примеры & Остатки, численность & Выручка, выпуск \\
		Формула средней & Хронологическая & Взвешенная арифм. \\
		Суммирование уровней & Недопустимо & Допустимо \\
		\bottomrule
	\end{tabular}
\end{center}

\textbf{Распространённая ошибка.} При расчёте средней хронологической нельзя
просто применить обычное арифметическое среднее: это занижает или завышает
результат, поскольку не учитывается, что крайние точки «представляют» лишь
половину прилегающего интервала.

% -----------------------------------------------------------
\subsection{Неравенство средних и выбор средней}
% -----------------------------------------------------------

Для строго положительных значений всегда выполняется:
\[
\bar{x}_{гарм} \;\leq\; \bar{x}_{геом} \;\leq\; \bar{x}_{арифм} \;\leq\; \bar{x}_{кв}
\]
Равенство достигается только при $x_1 = x_2 = \cdots = x_n$.

\bigskip
\begin{center}
	\renewcommand{\arraystretch}{1.4}
	\begin{tabular}{>{\bfseries}p{3.3cm} p{4.2cm} p{5cm}}
		\toprule
		Вид средней & Когда применять & Типичные задачи \\
		\midrule
		Арифметическая & Суммируемые признаки & Средний балл, ср. зарплата \\
		Взвешенная арифм. & Разные частоты/веса & Ср. цена по объёмам продаж \\
		Медиана & Скошенные распред., выбросы & Доходы населения, цены на жильё \\
		Мода & Наиболее частое значение & Самый популярный товар \\
		Геометрическая & Мультипликативный рост & Индексы, доходность \\
		Гармоническая & Фиксированный знаменатель & Ср. скорость, ср. цена \\
		\bottomrule
	\end{tabular}
\end{center}

\section{Меры положения: квантили и перцентили}
% -----------------------------------------------------------
\subsection{Квартили}
% -----------------------------------------------------------


	\textbf{Квартили} делят упорядоченный ряд на четыре равные части, по 25\% каждая.
	\begin{itemize}[noitemsep]
		\item $Q_1$ (нижний (первый) квартиль) - 25-й перцентиль.
		\item $Q_2$ (Второй квартиль) - 50-й перцентиль, он же медиана.
		\item $Q_3$ (верхний (третий) квартиль) - 75-й перцентиль.
	\end{itemize}


\subsubsection{Вычисление квартилей}

\textbf{В дискретном ряду}

\begin{enumerate}
	\item Упорядочить данные по возрастанию.
	\item Вычислить позицию квартиля: $L_i = i \cdot n / 4$, $\; i = 1,2,3$.
	\item Если $L_i$ целое: $Q_i = \bigl(x_{L_i} + x_{L_i+1}\bigr)/2$.\\
	Если $L_i$ дробное: $Q_i = x_{\lceil L_i \rceil}$.
\end{enumerate}


	\textbf{Пример.} Данные ($n=10$): 2, 5, 6, 8, 9, 11, 14, 17, 20, 25.
	\begin{align*}
		Q_1:&\quad L_1 = 2{,}5 \;\to\; Q_1 = x_3 = 6 \\
		Q_2:&\quad L_2 = 5{,}0 \;\to\; Q_2 = \tfrac{x_5+x_6}{2} = \tfrac{9+11}{2} = 10 \\
		Q_3:&\quad L_3 = 7{,}5 \;\to\; Q_3 = x_8 = 17
	\end{align*}

\textbf{В интервальном ряду}

\begin{enumerate}
	\item Вычислить интвервал искомого квартиля: $L_i = i \cdot n / 4$, $\; i = 1,2,3$.
	\item Найти первый интервал у которого $ S_i \geq  L_i$
	\item Используем формулу $ 	Q_i = x_0 + h \cdot \frac{L_i - S_{ПС-1}}{f_{ПС}}  $.
\end{enumerate}


% -----------------------------------------------------------
\subsection{Межквартильный размах (IQR)}
% -----------------------------------------------------------

\begin{equation}
	IQR = Q_3 - Q_1
\end{equation}

Показывает диапазон, в котором лежат центральные 50\% выборки.

\noindent\textbf{Правило выбросов Тьюки:}
\begin{align*}
	\text{Нижнее значение:} &\quad \text{НЗ} = Q_1 - 1{,}5 \cdot IQR \\
	\text{Верхнее значение:} &\quad \text{ВЗ} = Q_3 + 1{,}5 \cdot IQR
\end{align*}
Значения за пределами усов считаются выбросами.

% -----------------------------------------------------------
\subsection{Перцентили}
% -----------------------------------------------------------


	$P_k$ - $k$-й перцентиль: $k\%$ наблюдений находятся ниже этого значения.


\noindent Позиция $k$-го перцентиля:
\begin{equation}
	L_k = \frac{k \cdot n}{100}
\end{equation}

Применяется то же правило округления, что и для квартилей.


	\textbf{Интерпретация.} Если $P_{90} = 85$ баллов, то 90\% студентов набрали
	не более 85~баллов. Децили: $D_k = P_{10k}$. Квинтили: $V_k = P_{20k}$.


% =============================================================
\section{Box-Plot (ящик с усами, диаграмма Тьюки)}
% =============================================================

\subsection{Определение и элементы}


	\textbf{Box-plot}  - графическое
	представление распределения данных через пятичисловую сводку:
	\[
	\bigl\{\; x_{\min}^{*},\; Q_1,\; Q_2,\; Q_3,\; x_{\max}^{*} \;\bigr\}
	\]
	где $x_{\min}^{*}$ и $x_{\max}^{*}$ - значения усов.


\noindent\textbf{Элементы:}

\begin{enumerate}
	\item \textbf{Прямоугольник (ящик)}: от $Q_1$ до $Q_3$, ширина $= IQR$.
	\item \textbf{Линия внутри ящика}: медиана $Q_2$.
	\item \textbf{Усы}: линии от краёв ящика до ближайших наблюдений внутри
	значений Тьюки.
	\item \textbf{Выбросы}: значения за пределами усов - отдельные точки
	($\circ$ или $*$).
\end{enumerate}

\subsection{Алгоритм построения}

\begin{enumerate}
	\item Вычислить $Q_1$, $Q_2$, $Q_3$; затем $IQR = Q_3 - Q_1$.
	\item Рассчитать значения Тьюки: $\text{НЗ} = Q_1 - 1{,}5\cdot IQR$;
	$\;\text{ВЗ} = Q_3 + 1{,}5\cdot IQR$.
	\item Определить усы:
	\[
	\text{Нижний ус} = \min\{x_i \mid x_i \geq \text{НЗ}\}, \qquad
	\text{Верхний ус} = \max\{x_i \mid x_i \leq \text{ВЗ}\}.
	\]
	\item Отметить выбросы: все $x_i < \text{НЗ}$ или $x_i > \text{ВЗ}$.
	\item Нарисовать прямоугольник [$Q_1$;\,$Q_3$], линию $Me$, усы, выбросы.
\end{enumerate}

\subsection{Пример построения box-plot}


	\textbf{Данные} (ежемесячная выручка, тыс. р., $n=10$):
	\[
	15,\; 18,\; 20,\; 22,\; 25,\; 27,\; 29,\; 30,\; 35,\; 60
	\]
	\begin{align*}
		Q_1 &= x_3 = 20, \quad Q_2 = \tfrac{25+27}{2} = 26, \quad Q_3 = x_8 = 30\\
		IQR &= 30 - 20 = 10 \\
		\text{НЗ} &= 20 - 15 = 5, \qquad \text{ВЗ} = 30 + 15 = 45\\
		\text{Нижний ус} &= 15, \qquad \text{Верхний ус} = 35\\
		\text{Выброс:} &\quad 60 > 45
	\end{align*}
	
	\begin{center}
		\begin{tikzpicture}
			\begin{axis}[
				width=11cm, height=3.8cm,
				xmin=0, xmax=70, ymin=0, ymax=1,
				ytick=\empty, ylabel={},
				xlabel={Выручка, тыс. р.},
				axis lines=left,
				xtick={0,10,20,30,40,50,60,70}
				]
				\addplot[fill=blue!15, draw=black, line width=1pt]
				coordinates {(20,0.25)(30,0.25)(30,0.75)(20,0.75)(20,0.25)};
				\addplot[draw=red!80!black, line width=2pt]
				coordinates {(26,0.25)(26,0.75)};
				\addplot[draw=black, line width=1pt] coordinates {(15,0.5)(20,0.5)};
				\addplot[draw=black, line width=1pt] coordinates {(15,0.32)(15,0.68)};
				\addplot[draw=black, line width=1pt] coordinates {(30,0.5)(35,0.5)};
				\addplot[draw=black, line width=1pt] coordinates {(35,0.32)(35,0.68)};
				\addplot[mark=*, mark options={fill=red,scale=1.8}, draw=red]
				coordinates {(60,0.5)};
				\node[above] at (axis cs:60,0.72) {\footnotesize выброс};
				\node[above] at (axis cs:26,0.76) {\footnotesize $Q_2{=}26$};
				\node[below] at (axis cs:20,0.24) {\footnotesize $Q_1{=}20$};
				\node[below] at (axis cs:30,0.24) {\footnotesize $Q_3{=}30$};
				\node[below] at (axis cs:15,0.24) {\footnotesize 15};
				\node[below] at (axis cs:35,0.24) {\footnotesize 35};
			\end{axis}
		\end{tikzpicture}
	\end{center}


\subsection{Интерпретация box-plot}

\begin{center}
	\renewcommand{\arraystretch}{1.4}
	\begin{tabular}{p{5cm} p{7.5cm}}
		\toprule
		\textbf{Визуальный признак} & \textbf{Интерпретация} \\
		\midrule
		Медиана смещена к правому краю "ящика" & Правосторонняя (положительная) асимметрия \\
		Медиана смещена к правому краю "ящика" & Левосторонняя (отрицательная) асимметрия \\
		Медиана по центру ящика & Симметричное распределение \\
		Длинный правый ус & Хвост распределения направлен вправо \\
		Значения за усами & Выбросы, требующие проверки \\
		Большой IQR & Высокая вариация в «среднем» диапазоне \\
		Наложение ящиков при сравнении & Схожие распределения в группах \\
		\bottomrule
	\end{tabular}
\end{center}


	\textbf{Практическая рекомендация.} Сравнение нескольких ящиков на одном графике позволяет
	быстро сопоставить группы: регионы, периоды, продукты, по медиане, разбросу
	и наличию выбросов. При этом нет необходимости проверять нормальность
	распределения.
	
% =============================================================
\section{Статистические индексы}
% =============================================================


	\textbf{Индекс} - относительная величина, характеризующая соотношение двух
	уровней одного явления во времени, пространстве или по отношению к нормативу:
	\[
	I = \frac{\text{Текущий уровень}}{\text{Базовый уровень}}
	\]
	Результат $> 1$ означает рост, $< 1$ - снижение, $= 1$ - отсутствие изменений.


% -----------------------------------------------------------
\subsection{Простой, индивидуальный индекс}
% -----------------------------------------------------------

Для отдельного товара:

\begin{equation}
	i_p = \frac{p_t}{p_0}, \qquad i_q = \frac{q_t}{q_0}
\end{equation}

% -----------------------------------------------------------
\subsection{Базисный и цепной индексы}
% -----------------------------------------------------------

\noindent\textbf{Базисный индекс} - каждый период сравнивается с одним и тем же
базовым периодом~0:
\begin{equation}
	I^{\text{баз}}_t = \frac{y_t}{y_0}
\end{equation}

\noindent\textbf{Цепной индекс} - каждый период сравнивается с предыдущим:
\begin{equation}
	I^{\text{цеп}}_t = \frac{y_t}{y_{t-1}}
\end{equation}

\noindent\textbf{Взаимосвязь:}
\begin{equation}
	I^{\text{баз}}_t = \prod_{k=1}^{t} I^{\text{цеп}}_k
\end{equation}


	\textbf{Пример.} Динамика цен (база: 2020~г.):
	\begin{center}
		\begin{tabular}{ccccc}
			\toprule
			Год & Цена, р. & Баз. индекс & Цепной индекс & Темп роста, \% \\
			\midrule
			2020 & 100 & 1{,}000 & -   & -   \\
			2021 & 108 & 1{,}080 & 1{,}080 & 8{,}0 \\
			2022 & 112 & 1{,}120 & 1{,}037 & 3{,}7 \\
			2023 & 120 & 1{,}200 & 1{,}071 & 7{,}1 \\
			\bottomrule
		\end{tabular}
	\end{center}
	Проверка проверим результат: $1{,}080 \times 1{,}037 \times 1{,}071 \approx 1{,}200$.


% -----------------------------------------------------------
\subsection{Агрегатный индекс цен Ласпейреса}
% -----------------------------------------------------------


	\textbf{Индекс Ласпейреса} взвешивает цены текущего периода по объёмам
	\emph{базового} периода.


\begin{equation}
	I_P^L = \frac{\displaystyle\sum_{i} p_{it} \cdot q_{i0}}
	{\displaystyle\sum_{i} p_{i0} \cdot q_{i0}}
\end{equation}

\noindent\textbf{Экономический смысл:} во сколько раз изменилась бы стоимость
потребительской корзины \emph{базового} периода при текущих ценах.

\noindent\textbf{Свойства:}
\begin{itemize}[noitemsep]
	\item Корзина фиксирована на уровне базового периода.
	\item Завышает инфляцию, игнорируя эффект замещения.
	\item Применяется в ИПЦ (CPI) большинства стран.
\end{itemize}

% -----------------------------------------------------------
\subsection{Агрегатный индекс цен Пааше}
% -----------------------------------------------------------


	\textbf{Индекс Пааше} взвешивает цены по объёмам \emph{текущего} периода.


\begin{equation}
	I_P^P = \frac{\displaystyle\sum_{i} p_{it} \cdot q_{it}}
	{\displaystyle\sum_{i} p_{i0} \cdot q_{it}}
\end{equation}

\noindent\textbf{Экономический смысл:} во сколько раз изменилась стоимость
текущей корзины по сравнению с базовыми ценами.

\noindent\textbf{Сравнение индексов:}
\begin{center}
	\renewcommand{\arraystretch}{1.4}
	\begin{tabular}{p{3.5cm} p{4.5cm} p{4.5cm}}
		\toprule
		& \textbf{Ласпейрес} & \textbf{Пааше} \\
		\midrule
		Корзина & Базовый период & Текущий период \\
		Смещение & Завышает (при росте цен) & Занижает (при росте цен) \\
		Применение & CPI, ИПЦ & Дефлятор ВВП \\
		Данные & Нужны $q_{i0}$ & Нужны $q_{it}$ \\
		\bottomrule
	\end{tabular}
\end{center}

% -----------------------------------------------------------
\subsection{Индекс Фишера (идеальный)}
% -----------------------------------------------------------


	\textbf{Индекс Фишера} - среднее геометрическое индексов Ласпейреса и Пааше,
	устраняющее их систематические смещения.


\begin{equation}
	I_P^F = \sqrt{I_P^L \cdot I_P^P}
\end{equation}

Считается «идеальным», так как удовлетворяет тестам обратимости во времени
и цикличности.

% -----------------------------------------------------------
\subsection{Индекс физического объёма}
% -----------------------------------------------------------

\begin{equation}
	I_Q^L = \frac{\displaystyle\sum_{i} p_{i0} \cdot q_{it}}
	{\displaystyle\sum_{i} p_{i0} \cdot q_{i0}}
\end{equation}

Измеряет изменение физического объёма при ценах базового периода.

% -----------------------------------------------------------
\subsection{Взаимосвязь индексов товарооборота}
% -----------------------------------------------------------

\begin{equation}
	I_{PQ} = I_P^L \cdot I_Q^L
\end{equation}

где $\displaystyle I_{PQ} = \frac{\sum p_{it} q_{it}}{\sum p_{i0} q_{i0}}$ -
индекс товарооборота (в текущих ценах).


	\textbf{Пример.} Два товара, базовый и отчётный периоды:
	\begin{center}
		\begin{tabular}{ccccc}
			\toprule
			Товар & $p_0$ & $q_0$ & $p_t$ & $q_t$ \\
			\midrule
			A & 10 & 100 & 12 & 110 \\
			B & 20 &  50 & 22 &  45 \\
			\bottomrule
		\end{tabular}
	\end{center}
	\begin{align*}
		\sum p_0 q_0 &= 10 \times 100 + 20 \times 50 = 2000 \\
		\sum p_t q_0 &= 12 \times 100 + 22 \times 50 = 2300 \\
		\sum p_0 q_t &= 10 \times 110 + 20 \times 45 = 2000 \\
		\sum p_t q_t &= 12 \times 110 + 22 \times 45 = 2310 \\[6pt]
		I_P^L &= \frac{2300}{2000} = 1{,}15 \;\; (\text{рост цен на } 15\%) \\
		I_Q^L &= \frac{2000}{2000} = 1{,}00 \;\; (\text{объём не изменился}) \\
		I_{PQ} &= \frac{2310}{2000} = 1{,}155 \\
		I_P^P &= \frac{2310}{2000} = 1{,}155 \\
		I_P^F &= \sqrt{1{,}15 \times 1{,}155} \approx 1{,}152
	\end{align*}


% -----------------------------------------------------------
\subsection{Индексы среднего уровня: фиксированного, переменного состава
	и структурных сдвигов}
% -----------------------------------------------------------

\subsubsection{Индекс переменного состава}


	\textbf{Индекс переменного состава} характеризует изменение среднего уровня
	с учётом одновременного изменения \emph{и} значений показателя, \emph{и}
	структуры совокупности.


\begin{equation}
	I_{\text{пер}} = \frac{\bar{x}_t}{\bar{x}_0}
	= \frac{\displaystyle\sum_i x_{it} \cdot d_{it}}
	{\displaystyle\sum_i x_{i0} \cdot d_{i0}}
\end{equation}

где $d_{it}$, $d_{i0}$ - доли $i$-й группы в текущем и базовом периодах.

\subsubsection{Индекс фиксированного состава}


	\textbf{Индекс фиксированного состава} измеряет изменение среднего уровня
	только за счёт изменения значений показателя при \emph{неизменной} структуре
	базового периода.


\begin{equation}
	I_{\text{фикс}} =
	\frac{\displaystyle\sum_i x_{it} \cdot d_{i0}}
	{\displaystyle\sum_i x_{i0} \cdot d_{i0}}
\end{equation}

\subsubsection{Индекс структурных сдвигов}


	\textbf{Индекс структурных сдвигов} показывает, как изменился бы средний
	уровень только за счёт изменения структуры при \emph{неизменных} значениях
	показателя базового периода.


\begin{equation}
	I_{\text{стр}} =
	\frac{\displaystyle\sum x_{0} \cdot d_{t}}
	{\displaystyle\sum_i x_{0} \cdot d_{0}}
\end{equation}

\subsubsection{Взаимосвязь трёх индексов}

\begin{equation}
	\boxed{I_{\text{пер}} = I_{\text{фикс}} \cdot I_{\text{стр}}}
\end{equation}


	\textbf{Пример.} Средняя заработная плата по двум подразделениям:
	
	\begin{center}
		\renewcommand{\arraystretch}{1.3}
		\begin{tabular}{ccccc}
			\toprule
			Подразделение & $x_{i0}$, тыс. р. & $d_{i0}$ & $x_{it}$, тыс. р. & $d_{it}$ \\
			\midrule
			I  & 50 & 0{,}40 & 55 & 0{,}60 \\
			II & 80 & 0{,}60 & 85 & 0{,}40 \\
			\bottomrule
		\end{tabular}
	\end{center}
	
	\begin{align*}
		\bar{x}_0 &= 50\times0{,}40 + 80\times0{,}60 = 20+48 = 68{,}0 \text{ тыс. р.}\\
		\bar{x}_t &= 55\times0{,}60 + 85\times0{,}40 = 33+34 = 67{,}0 \text{ тыс. р.}\\[4pt]
		I_{\text{пер}} &= \frac{67{,}0}{68{,}0} \approx 0{,}985 \quad (-1{,}5\%)\\[4pt]
		I_{\text{фикс}} &= \frac{55\times0{,}40+85\times0{,}60}{68{,}0}
		= \frac{22+51}{68} = \frac{73}{68} \approx 1{,}074 \quad (+7{,}4\%)\\[4pt]
		I_{\text{стр}} &= \frac{50\times0{,}60+80\times0{,}40}{68{,}0}
		= \frac{30+32}{68} = \frac{62}{68} \approx 0{,}912 \quad (-8{,}8\%)\\[4pt]
		\text{Проверка:} &\quad 1{,}074 \times 0{,}912 \approx 0{,}979 \approx 0{,}985
		\;\; 
	\end{align*}
	
	\textbf{Вывод.} Зарплата в каждом подразделении выросла на 7{,}4\% - индекс
	фиксированного состава. Однако сокращение доли высокооплачиваемого
	подразделения~II снизило среднюю зарплату на 8{,}8\% - структурный сдвиг,
	что дало итоговое снижение средней зарплаты по предприятию на 1{,}5\%.


\subsection{Интерпретация индексов состава}

\begin{center}
	\renewcommand{\arraystretch}{1.4}
	\begin{tabular}{p{4.2cm} p{8cm}}
		\toprule
		\textbf{Индекс} & \textbf{Интерпретация} \\
		\midrule
		$I_{\text{пер}} > 1$ & Средний уровень в целом вырос \\
		$I_{\text{пер}} < 1$ & Средний уровень в целом снизился \\
		$I_{\text{фикс}} > 1$, $I_{\text{стр}} < 1$ & Рост показателя внутри групп
		компенсируется неблагоприятным структурным сдвигом \\
		$I_{\text{фикс}} < 1$, $I_{\text{стр}} > 1$ & Структурные сдвиги улучшают
		средний уровень, несмотря на падение показателя \\
		$I_{\text{стр}} = 1$ & Структура не изменилась \\
		\bottomrule
	\end{tabular}
\end{center}

% =============================================================
\section{ИПЦ и дефлятор ВВП как приложение индексной теории}
% =============================================================

\subsection{Индекс потребительских цен}

\begin{equation}
	\text{CPI}_t = \frac{\displaystyle\sum_i p_{it}\, q_{i0}}
	{\displaystyle\sum_i p_{i0}\, q_{i0}} \times 100
\end{equation}

По форме - индекс Ласпейреса. Используется для измерения инфляции,
индексации зарплат и дефлятирования временных рядов.

\noindent\textbf{Дефлятирование:}
\begin{equation}
	\text{Реальная величина}_t = \frac{\text{Номинальная величина}_t}{\text{CPI}_t / 100}
\end{equation}

\subsection{Дефлятор ВВП}

\begin{equation}
	\text{Дефлятор ВВП}_t
	= \frac{\displaystyle\sum_i p_{it}\, q_{it}}{\displaystyle\sum_i p_{i0}\, q_{it}} \times 100
\end{equation}

По форме - индекс Пааше. Охватывает все товары и услуги, входящие в ВВП
(шире потребительской корзины CPI).

% =============================================================
\section{Таблица формул}
% =============================================================

\begin{center}
	\renewcommand{\arraystretch}{1.8}
	\begin{tabular}{p{4.5cm} p{8cm}}
		\toprule
		\textbf{Показатель} & \textbf{Формула} \\
		\midrule
		Ср. арифметическое & $\bar{x} = \dfrac{\sum x_i}{n}$ \\
		Взвешенное ср. арифм. & $\bar{x}_w = \dfrac{\sum x_i f_i}{\sum f_i}$ \\
		Медиана (сгр.) &
		$\text{Me} = L + \dfrac{n/2 - F_{m-1}}{f_m}\cdot h$ \\
		Мода (сгр.) &
		$\text{Mo} = L + \dfrac{f_m - f_{m-1}}{(f_m-f_{m-1})+(f_m-f_{m+1})}\cdot h$ \\
		Ср. геометрическое & $\bar{x}_{\text{геом}}  = \bigl(\prod x_i\bigr)^{1/n}$ \\
		Ср. гармоническое & $\bar{x}_{\text{гарм}}  = n\,/\,\sum(1/x_i)$ \\
		Ср. хронологическая (равн.) &
		$\bar{x}_{\text{хр}} = \dfrac{x_1/2 + x_2 + \cdots + x_{n-1} + x_n/2}{n-1}$ \\
		Ср. хронологическая (неравн.) &
		$\bar{x}_{\text{хр}} = \dfrac{\sum \frac{x_j+x_{j+1}}{2}\cdot t_j}{\sum t_j}$ \\
		Общая формула порядковой статистики & $ \text{ПС}_i = x_0 + h \cdot \frac{\frac{i}{j}  \cdot n - S_{ПС-1}}{f_{ПС}} $ \\
		Квартиль & $Q_i = i\cdot \frac{n}{ 4}$ \\
		Перцентиль & $P_k = k\cdot \frac{n}{100}$ \\
		IQR & $IQR = Q_3 - Q_1$ \\
		"Усы" по Тьюки & $Q_1 \pm 1{,}5\cdot IQR$ \\
		Индекс Ласпейреса & $I_P^L = \dfrac{\sum p_t q_0}{\sum p_0 q_0}$ \\
		Индекс Пааше & $I_P^P = \dfrac{\sum p_t q_t}{\sum p_0 q_t}$ \\
		Индекс Фишера & $I_P^F = \sqrt{I_P^L \cdot I_P^P}$ \\
		Инд. перем. состава & $I_{\text{пер}} = \dfrac{\sum x_t d_t}{\sum x_0 d_0}$ \\
		Инд. фикс. состава & $I_{\text{фикс}} = \dfrac{\sum x_t d_0}{\sum x_0 d_0}$ \\
		Инд. структ. сдвигов & $I_{\text{стр}} = \dfrac{\sum x_0 d_t}{\sum x_0 d_0}$ \\
		Связь трёх инд. & $I_{\text{пер}} = I_{\text{фикс}} \cdot I_{\text{стр}}$ \\
		\bottomrule
	\end{tabular}
\end{center}




\end{document}
